% Technical report outline -- proposed structure, not a completed manuscript.
% Author-confirmed emphasis: full experimental history and decision rationale.
% This is a section fragment for the existing Report/main.tex template.
% Chinese comments are writing notes; they introduce no LaTeX dependencies.
% Prepared from repository evidence available on 2026-10-21 (Asia/Taipei).
%
% Proposed title:
% Training Llama 3.2 1B from Scratch: An Experimental Study of Optimization,
% Token Efficiency, and Compute Cost
%
% 已確認的敘事主線：從隨機初始化開始，記錄每輪改善 held-out perplexity
% 的問題、假說、對照、結果與下一步決策。依實驗階段組織，保留負面結果。
% 正文呈現影響後續決策的實驗；完整時間線、所有候選與成本放附錄。
%
% Abstract writing plan (write after results are reconciled):
% 問題與限制 -> 逐階段研究方式 -> 主要結果 -> token/compute 取捨 -> 限制。
% 不將最終配方的多項變更寫成某個單一方法的因果增益。

\section{Introduction}
\label{sec:introduction}
% 動機：在固定 Llama 架構、8192 context 與有限訓練 tokens 下改善泛化。
% 說明隨機初始化學生、預訓練教師與 checkpoint continuation 的區別。
% 貢獻定位為可追溯的實驗研究、決策分析與可重現配方；不預設演算法新穎性。
\subsection{Research Questions and Objectives}
% 哪些變更改善相同 tokens 的 NLL？哪些降低成本？短程排序能否轉移到長程？
\subsection{Overview of the Experimental Progression}
% 用一張實驗演進圖標示各階段、關鍵分支與採用／停止的理由。

\section{Experimental Setup and Evaluation Protocol}
\label{sec:setup}
\subsection{Model, Tokenizer, and Random Initialization}
% 記錄實際參數量、架構、tied embedding/head、tokenizer revision、seed 與8192 context。
% 初始化分布的消融放後續實驗章；此處先交代共同設定。
\subsection{Canonical Data Split and Evaluation Subsets}
% 完整 train 文件先 canonical shuffle(seed42)，最後50000篇保留，再 tokenize/pack。
% 分清固定1k、額外9k、合併10k與full50k；列出各階段使用情形與重用情形。
\subsection{Metrics and Comparison Criteria}
% 主指標為 sum(NLL)/sum(valid predictions)，PPL取其指數。
% 分開比較相同token endpoint、相同compute budget、完整allocation成本與step throughput。
% 區分 input tokens、loss-bearing prediction tokens、teacher tokens 與eval tokens。
\subsection{Experiment Categories and Resource Accounting}
% 區分fresh run、pilot、共享checkpoint分支及明確授權的continuation。
% 說明原始預算與後續擴充；記錄實際rounded endpoints、失敗allocation及總搜尋成本。
% 不把超過原始64H的例外run描述為達成64H目標。

\section{Data Preparation and Early Baselines}
\label{sec:baseline}
\subsection{Data Integrity, Tokenization, and Packing}
% 精簡交代canonical ordering、損壞shard修復、split驗證、BOS/EOS及padding/loss mask。
% 將具體故障、hash與恢復時間線移至附錄；正文說明對實驗有效性的影響。
\subsection{AdamW Baseline and Initial Optimizer--Scheduler Screening}
% 從最早AdamW基線，接到Muon及cosine/linear/WSD比較。
% 報告匹配token endpoint、搜尋範圍、seed與評估口徑；不直接比較不等長短跑。
\subsection{Raw versus Filtered Training Data}
% 說明DCLM篩選動機、實際保留規則、來源組成與raw/filtered對照。
% 僅納入確實完成的比較；raw-filtered-3b的filtered arm已取消，不能當完整雙臂結果。
% 討論來源分布與token availability，避免概括成「品質篩選普遍無效」。

\section{Improving Training Throughput}
\label{sec:systems}
\subsection{Profiling and Implementation Changes}
% 沿時間線介紹資料reader/prefetch、microbatch/accumulation、compile、attention與optimizer實作。
% 逐項確認變更是否保留目標、mask、資料順序與梯度語義。
\subsection{Correctness Checks and Measured Speedups}
% 以完整update、冷啟動、峰值記憶體及allocation成本呈現，保持計時口徑一致。
% performance-ab-20261010/results.md可作早期短程測量，不能外推為6B品質等價證明。
\subsection{How Efficiency Changed the Experiment Plan}
% 說明節省的成本如何支援後續候選數／訓練長度；實測與投影分列。

\section{Learning-Rate Schedules and Training-Horizon Transfer}
\label{sec:schedules}
\subsection{WSD, Cooldown Shape, and Phase Splitting}
% 說明各輪tail比例、plateau、第二phase/reheating，以及觀察如何改變下一輪設計。
\subsection{Learning Rate, Gradient Clipping, and Initialization Studies}
% 依實際先後交代LR/clipping研究與初始化配方sweep，涵蓋residual-half與組合確認。
% 單因素改善、組合結果與seed43確認分開呈現。
\subsection{Proxy Failure and Revised Horizon Calibration}
% 保留stable proxy與完整WSD不一致的經驗，交代為何重做bracket與true-horizon校準。
% 區分實測、擬合及1--6B外推；fit-window敏感度不可稱統計信賴區間。

\section{Batch Policy, Attention Masking, and Optimizer Refinement}
\label{sec:refinement}
\subsection{Batch-Size Transitions and Schedule Clocks}
% 固定batch對照與BS16->64的切換時機；比較固定／縮放LR。
% 明確區分token-clock與update-clock；batch與decay同時變動時承認混淆因素。
\subsection{Intra-Document Attention}
% 從document_spans定義block-diagonal causal mask，不從token ID推測邊界。
% 展示matched-token對照、文件長度分桶、seed43確認及3B/4B轉移證據。
\subsection{Muon, NorMuon, and Auxiliary AdamW Learning Rates}
% 交代矩陣與embedding/head/norm分組、NS5、normalization與auxiliary AdamW LR x2。
% 將PE-only/NorMuon-only和組合比較分開，不把複合變更歸因給單一項。
\subsection{Weight Averaging, Weight Decay, and Negative Findings}
% 納入EMA/有限窗口平均及WD相關對照，呈現未採用的原因與適用範圍。
\subsection{Scaling the Selected Recipe to 3B, 4B, and 6B Tokens}
% 完整交代每次fresh初始化、LR、batch、phase與實際tokens；3B/4B/6B是訓練量。
% 相同token的4B AdamW x2與step-tail30可作重點；跨token預算曲線僅描述整體演進。

\section{Further Method Exploration and Knowledge Distillation}
\label{sec:distillation}
\subsection{Alternative Optimizers and Training Objectives}
% 納入Dion3、AdaMuon、SOAP、schedule-free與CCE等實際研究。
% 完成1B、提早停止、程式故障與未執行提案分列；不以不同endpoint硬排名。
\subsection{From the Initial Teacher to the Official 3B Teacher}
% 說明教師選擇改變的動機、學生隨機初始化、CE+KL目標、temperature與teacher權重。
% 教師預訓練屬於外部資源，不能描述成整個訓練流程只使用random weights。
\subsection{Producer--Consumer Transport and Single-Allocation KD}
% 交代8+8 hidden-state transport、pipeline限制及後續R5融合CE/KL的設計決策。
% 品質與工程可靠性分開討論；含失敗成本及師生各自處理tokens。
\subsection{The 6B KD Result and Its Cost}
% 已核實R5：6000082944 student inputs；combined10k PPL=12.3989218778246。
% 相同10k上的non-KD6B=12.9078477387367；PPL相對降低約3.94%。
% R5實際64.711111H，屬72H授權；不是原64H目標內完成。
% 沒有matched unaccelerated6B KD reference，不能據此宣稱加速完全無損。
\subsection{Adaptive Critical Batch and Token-Bound Early Decay}
% 交代C0 shadow/C1 half/C2 one及後續6B C2 early-decay的假說與設計。
% First decay綁2399928320累積nonpadding input tokens，不綁W&B/local updates。
% 目前已讀到的C2檔案僅startup-status；正文數字需另核實最終checkpoint/eval/accounting。

\section{Consolidated Results and Decision Analysis}
\label{sec:results}
\subsection{Milestones in the Training Recipe}
% 一張里程碑表：日期、parent/control、變更、tokens、eval subset、NLL/PPL、allocation成本與決策。
% 另列formal identities與分支關係，避免將共享prefix當獨立重跑。
\subsection{Token Efficiency and Compute Efficiency}
% 在同一eval口徑畫PPL vs student inputs與PPL vs H200-hours。
% KD顯式加入teacher allocation；分開單次run成本、搜尋成本與歷史teacher成本。
\subsection{Which Findings Changed Subsequent Decisions?}
% 選3--5個有控制證據的結論，逐一連結假說、觀察與採用／停止理由。
% 加入短程排序轉移、組合非加成、mask與batch/schedule交互作用的限制。

\section{Discussion and Limitations}
\label{sec:discussion}
\subsection{What Generalizes Beyond This Experiment?}
% 將可重用的方法論與本資料／token horizon下的觀察區分。
\subsection{Statistical and Experimental Limitations}
% 大多single-seed、有限seed43確認、反覆重用eval subsets與自適應選擇。
% 文件bootstrap不能取代training seed變異；缺少控制格就明說。
% PPL不等同downstream能力；外部預訓練teacher與資料重疊風險不能略去。
\subsection{Operational Failures and Reproducibility}
% 精選會改變成本或結果可信度的案例；其餘job細節放附錄。

\section{Conclusion}
\label{sec:conclusion}
% 回答引言研究問題：哪些實驗有效、哪些未支持假說、最終配方的品質／成本。
% 僅以已核實終點總結，不將尚未完成或未核實的候選描述為改善。

% Related work plan:
% 將背景按pretraining/optimizer/schedule/data selection/KD分主題引用，
% 放各章開頭或整合為短獨立章。正式成稿時逐筆查證文獻與BibTeX。
% 現有references.bib是template示例，不是本研究的已驗證參考文獻。

\appendix
\section{Complete Experiment Timeline and Run Registry}
% 所有formal/pilot/branch/failed/cancelled identities、日期、W&B與artifact位置。
\section{Resolved Configurations and Hyperparameter Search Spaces}
% 對照表、初始化、optimizer routing、scheduler與batch planner公式／pseudocode。
\section{Data and Evaluation Provenance}
% pinned revisions、split/index hashes、packing/document spans、subset與NLL聚合方式。
\section{Additional Ablations and Learning Curves}
% 全部候選、seed43、文件長度／來源分析與未採用方法。
\section{Resource Ledger, Failures, and Recovery Records}
% 實際Slurm allocation、預算修訂、失敗／取消成本、teacher與student分帳。
\section{Checkpoint, Export, and Reproduction Details}
% 完整state保存、HF reload驗證、source/config seals與可重現指令。

% Priority figures/tables:
% F1: 實驗演進與決策圖。
% F2: 同eval下PPL vs cumulative student input tokens。
% F3: LR與batch對tokens的曲線，標註phase/decay/batch事件。
% F4: 關鍵matched-token消融與有限seed replication。
% F5: PPL vs allocation H200-hours，KD包含teacher成本。
% T1: 共同實驗設定與評估口徑。
% T2: 配方演進里程碑與採用／停止決策。
% T3: 最終non-KD/KD及已核實adaptive-batch結果。
